\chapter{Практика}
\label{ch:chap2}

% \begin{figure}[H]
%     \centering
%     \includegraphics{carrier.png}
%     \caption{Несущие колебания}
% \end{figure}

\section{Центральная предельная теорема}

Центральная предельная теорема (ЦПТ) — фундаментальное положение теории вероятностей, утверждающее, что сумма 
(или среднее значение) большого количества независимых, одинаково распределенных случайных величин имеет распределение, 
близкое к нормальному. \\

Для проверки теоремы сгенерируем выборки из равномерного распределения (реализации) и просуммируем их. Если теорема верна,
то получим случайную величину с распределением близким к нормальному:

\begin{figure}[H]
    \centering
    \includegraphics{1.png}
    \caption{Плотность распределения}
\end{figure}

Получили плотность распределения, близкую к нормальной, что подтверждает ЦПТ.

\section{АКФ по множеству реализаций}

Если есть множество реализаций СП, то АКФ можно вычислить по следующему алгоритму:

\begin{enumerate}
    \item Если среднее СП не ноль, то вычесть его из реализации.
    \item Посчитать АКФ для каждой реализации как $\sum_{n=0}^N x[n]*x[n+lag]$ и усреднить по времени, т.е $R = \frac{1}{N-lag}\sum_{n=0}^N x[n]*x[n+lag]$.
    \item Сложить полученные АКФ и усреднить по ансамблю, т.е необходимо разделить сумму на кол-во реализаций.
\end{enumerate}

\begin{figure}[H]
    \centering
    \includegraphics{2.png}
    \caption{АКФ}
\end{figure}

Видим один пик при нулевом лаге, что логично, ведь тогда процесс на 100\% совпадает сам с собой, и значения хорошо коррелируют.
При остальных лагах АКФ принимает значения близкие к нулю. Это говорит о том, что процесс не коррелирован с собой во времени, т.е 
значения не зависят от предыдущих. У процесса нет памяти.

\section{АКФ после свертки}

Канал можно рассматривать как систему, имеющую импульсную характеристику $h(t)$. Выполним свертку СВ с нормальным распределением
с какой-то ИХ и построим АКФ для полученной СВ

\begin{figure}[H]
    \centering
    \includegraphics{3.png}
    \caption{АКФ}
\end{figure}

На нижнем графике видим, что АКФ при лаге не равном нулю не принимает значения около нуля, т.е наш ипровизированный канал внес 
зависимость между значениями СВ.

\section{Интервал корреляции}

Интервал корреляции - интервал, на котором СВ сохраняет зависимость между своими значениями. Чем больше значение, тем процесс
меняется плавнее (соседние значение больше зависимы), чем меньше значения, тем процесс меняется хаотичнее. \\

Интервал корреляции можно вычислить как

\begin{figure}[H]
    \centering
    \includegraphics{4.png}
    \caption{Интервал корреляции}
\end{figure}

Посчитаем интервал корреляции для СВ с нормальным распределением:

\begin{figure}[H]
    \centering
    \includegraphics{5.png}
    \caption{Интервал корреляции}
\end{figure}

Интервал достаточно короткий, т.е корреляция между значениями СВ слабая и СВ ведет себя хаотично, что очень похоже на нормальное распределение.

\endinput 